% Generated draft: requires mathematical and source review.
\documentclass[11pt]{article}
\usepackage[margin=25mm]{geometry}
\usepackage{fontspec}
\setmainfont{Latin Modern Roman}
\usepackage{amsmath,amssymb,mathtools,unicode-math}
\setmathfont{Latin Modern Math}
\usepackage{enumitem,needspace,xurl,hyperref}
\setlength{\parindent}{0pt}
\setlength{\parskip}{5pt}
\setlength{\emergencystretch}{3em}
\newcommand{\N}{\mathbb{N}}
\newcommand{\Z}{\mathbb{Z}}
\newcommand{\Q}{\mathbb{Q}}
\newcommand{\R}{\mathbb{R}}
\newcommand{\C}{\mathbb{C}}
\newcommand{\F}{\mathbb{F}}
\newcommand{\A}{\mathbb{A}}
\newcommand{\PP}{\mathbb{P}}
\newcommand{\E}{\mathbb{E}}
\newcommand{\eps}{\varepsilon}
\newcommand{\dd}{\,\mathrm{d}}
\begin{document}
\section*{UnsolvedMath ID: 800207}
\textbf{Source problem code:} catalogue entry
\subsection{Definitions and mathematical statement}
Let \(M\) and \(N\) be strongly pseudoconvex CR manifolds of real dimension \(3\). Their CR structures consist of real rank-two distributions \(H_M\subset TM\) and \(H_N\subset TN\), with complex structures \(J_M\) and \(J_N\) on these distributions. Strong pseudoconvexity makes the horizontal distributions contact distributions. A smooth map \(f:M\to N\) is CR if, for every \(p\in M\),
\[
\dd f_p(H_{M,p})\subseteq H_{N,f(p)},\qquad
\dd f_p\circ J_{M,p}=J_{N,f(p)}\circ\dd f_p\quad\text{on }H_{M,p}.
\]
Given a smooth positive definite inner product \(g\) on a horizontal distribution \(H\), its Carnot--Carath\'eodory distance is locally defined by
\[
d_g(x,y)=\inf_{\gamma}\int_0^1\sqrt{g_{\gamma(t)}(\dot\gamma(t),\dot\gamma(t))}\,\dd t,
\]
where the infimum runs over absolutely continuous curves joining \(x\) to \(y\) with \(\dot\gamma(t)\in H_{\gamma(t)}\) almost everywhere. Distances may be considered separately on connected components. For CR-compatible choices one requires \(g(JX,JY)=g(X,Y)\).

\textbf{Source question.} Suppose \(f:M\to N\) is a homeomorphism that maps infinitesimal spheres for Carnot--Carath\'eodory metrics to infinitesimal spheres, as required by the source's definition of conformality. Must \(f\) be CR?

The phrase ``maps infinitesimal spheres to infinitesimal spheres'' is not given a limiting definition in the supplied record. One possible precise metric interpretation, after fixing distances \(d_M,d_N\), is the following. For sufficiently small \(r>0\), put
\[
L_f(p,r)=\sup_{d_M(p,q)=r}d_N(f(p),f(q)),\qquad
\ell_f(p,r)=\inf_{d_M(p,q)=r}d_N(f(p),f(q)).
\]
Metric conformality could then mean
\[
\limsup_{r\downarrow0}\frac{L_f(p,r)}{\ell_f(p,r)}=1
\qquad\text{for every }p\in M.
\]
Under this interpretation, a precise version asks whether every homeomorphism satisfying this condition is smooth and satisfies the CR identities above. This is a conditional formalization: the supplied source does not explicitly select this distortion condition, specify the horizontal metrics or their quantification, or state the regularity intended by the conclusion. It also imposes no explicit orientation condition and does not discuss whether conjugate-CR maps are to be included. These qualifications must be resolved before treating this as a uniquely determined statement.
\subsection{Short English statement}
The question attributed to Puqi Tang asks whether a conformal homeomorphism between strongly pseudoconvex CR three-manifolds must be CR, where conformality means preserving infinitesimal Carnot--Carath\'eodory spheres. The exact metric definition, metric choices, orientation conventions, and intended regularity remain unspecified. The registry labels the problem open; that label is supplied source data, not an independently verified status assessment.
\subsection{Sources}
\begin{itemize}
\item Supplied local UnsolvedMath registry record: permanent ID \(800207\), exact source code \texttt{catalogue entry}, title ``Several Complex Variables Problem List --- Problem 3,'' in the AMR Open Problem Lists collection. Supplied status: open; supplied proposed year: 1995.
\item Buzzard, Fornaess, and Krantz, \emph{The Several Complex Variables Problem List} (1995), Problem 3. Supplied source URL: \url{https://arxiv.org/abs/math/9509202}. The record reports extraction by \texttt{arxiv\_latex} and access on 2026-10-16; these are supplied provenance statements, not independently verified access or literature assessment.
\item The source item is explicitly attributed to Puqi Tang, \texttt{tang@math.purdue.edu}. Although the registry's proposer field lists Buzzard, Fornaess, and Krantz, the item itself credits Tang. Its introductory qualifications are ``(New),'' ``(See Section 11.3),'' and ``Liouville type assertion on CR 3-manifolds.'' The referenced Section 11.3 is not included in the supplied record.
\end{itemize}
\end{document}
